\chapter*{序言：当答案变得便宜}
\addcontentsline{toc}{chapter}{序言：当答案变得便宜}

2026 年 7 月 1 日，一门桥梁工程课程设计课进行结题汇报。学生展示的不再只是计算书和图纸：有桥梁拆除游戏，有影响线动画，有交互式计算工具，也有通过 AI 连接建模软件完成的桥梁模型。

作品很丰富，问题也同样具体。一个小组发现，AI 生成的弯矩图画错了；公式被放在不合适的位置；影响线的负值没有参与计算；继续修改时，一个失败版本还覆盖了接近完成的版本。讲到最后，学生说：

\begin{shiyu-inline}
``而且过程中就发现，你就真的得是自己会才行，要不然你也不知道 AI 做得对不对。像那个弯矩图。''

\hfill ——课程结题汇报，2026 年 7 月 1 日 00:25:20
\end{shiyu-inline}

这句话比“AI 会不会取代教师”更接近本书要处理的问题。

生成式 AI 已经可以迅速给出文本、代码、图片和模型。于是，\textbf{交出一个看起来完整的成品}不再足以证明学生经历了学习。反过来，禁止学生使用 AI 也不能恢复旧秩序：学生毕业后面对的正是开放工具、真实任务、多人协作和不断变化的工作流。

真正需要重做的是课堂里的学习回路。

本书所说的学习回路，是一个朴素的过程：\textbf{学生先尝试，系统或他人给出反馈，学生据此修改，然后解释自己为什么这样改。}在桥梁工程里，作品可以是一座游戏中的桥、一张弯矩图、一个有限元模型、一段施工动画或一份方案比较。作品不是终点，而是让思考变得可见的介质。

我过去的课堂更接近常见的工科教学：知识按章节展开，公式在黑板上推导，学生在作业和考试里复现。这样的教学能传递系统知识，却不保证知识进入工程直觉。2026 年 5 月的一次分享中，我谈到一个不太体面的观察：一些已经学过材料力学、结构力学和理论力学的学生，第一次打开搭桥游戏时，仍主要依靠“加一根杆试试”的方式探索。这个观察不能证明理论课无效。它提醒我：\textbf{会在纸面上求解，不等于会在一个开放问题中识别什么值得算。}

游戏、可视化和 AI 因此进入课堂。它们不是为了让课程显得新，也不是给公式加一层糖衣。它们分别解决三类问题：游戏压缩尝试与反馈之间的时间；可视化让不可见的关系可以操作；AI 降低制作原型的门槛，让更多学生能把想法变成可检查的作品。

这些工具也会制造新的错觉。画面漂亮，不等于模型正确；学生参与热闹，不等于理解发生；AI 能运行一段代码，不等于代码表达了正确的工程假设。关于数字游戏的研究提示，游戏对学习的平均作用为正，但效果取决于具体设计，而不是“用了游戏”这件事本身\cite{clark2016games}。关于反馈的研究也一再指出，反馈的类型、时机和学生是否真正使用反馈，比“给了多少评语”更重要\cite{shute2008feedback,carless2018literacy}。

所以，这不是一本成功案例集。书中会保留没做出来的自动桥型、画错的弯矩图、无法解释的模态、失控的版本和教师并不满意的作业。失败不是为了制造戏剧，而是因为课堂设计只有在这些地方才显出结构。

\section*{这本书写给谁}

第一读者是高校工科教师。你有自己的专业课，知道什么不能讲错；你可能已经试过几个 AI 工具，也可能正被“全面拥抱 AI”和“严禁 AI 代写”两种声音同时拉扯。你需要的不是更多工具名，而是一套能用于下一节课的设计方法。

教师发展工作者、专业负责人和研究生助教也可以把本书当作磨课材料。桥梁工程是贯穿全书的现场，但书中的任务、反馈、证据链和评价方法可以迁移到机械、建筑、交通、物理实验以及其他强调模型与责任的课程。

\section*{读完以后，请做一件事}

不要重做整门课。选一个学生经常“会算但不会判断”的知识点，把它改成一个 45--90 分钟的小任务。要求学生留下初次尝试、得到的反馈、一次修改和一段解释。允许使用 AI，但要求标明它参与了哪一步、学生怎样核验。

如果这个单元能让你看见学生原本藏在最终答案后面的思考，本书就完成了它最重要的工作。

桥在这里仍然是一个好隐喻，但我会谨慎使用它。桥梁把分开的地方建立成关系；课堂设计也把知识、行动、反馈和责任连在一起。隐喻不能替代证据，正如漂亮的渲染不能替代计算。它只负责提醒我们：真正重要的不是把内容从这一岸搬到另一岸，而是让学生能够自己走过去。

\begin{flushright}
ColdCat（大连交通大学）\\
2026 年 7 月，于大连
\end{flushright}
